\documentclass[11pt]{amsart}

\usepackage{amsmath,amssymb,amsthm,mathtools}
\usepackage[T1]{fontenc}
\usepackage{lmodern}
\usepackage[expansion=false]{microtype}
\setlength{\emergencystretch}{1em}
\usepackage[hidelinks]{hyperref}
\hypersetup{pdftitle={Moh's P3 curve is a global set-theoretic complete intersection},pdfauthor={}}

\newtheorem{theorem}{Theorem}[section]
\newtheorem{lemma}[theorem]{Lemma}
\newtheorem{corollary}[theorem]{Corollary}
\newtheorem{proposition}[theorem]{Proposition}
\theoremstyle{remark}
\newtheorem{remark}[theorem]{Remark}

\DeclareMathOperator{\ara}{ara}
\DeclareMathOperator{\adj}{adj}

\title{Moh's $P_3$ curve is a global set-theoretic complete intersection}
\date{Version 2 --- September 29, 2026}

\subjclass[2020]{Primary 13C40, 14M10; Secondary 13H05, 13D02, 13P10}
\keywords{set-theoretic complete intersection, arithmetic rank, Moh prime, Hilbert--Burch theorem, partly symmetric matrix, determinantal ideal}

\begin{document}

\begin{abstract}
Let $k$ be a field of characteristic zero and let $\mathfrak p$ be the kernel of the homomorphism $k[x,y,z]\to k[t]$ given by
\[
 (x,y,z)\longmapsto(t^6+t^{31},t^8,t^{10}).
\]
This note gives polynomials $H_1,H_2\in\mathbb Z[x,y,z]$, each of total degree $16$, such that
\[
 \mathfrak p^2\subseteq(H_1,H_2)\subseteq\mathfrak p.
\]
Thus the affine curve is a set-theoretic complete intersection and $\ara(\mathfrak p)=2$. An explicit invertible polynomial row operation places a Hilbert--Burch matrix in the partly symmetric setting of Valla and Kleiman--Ulrich. Their determinantal construction then gives the square containment. The same equations define the corresponding formal prime up to radical, addressing the question for this curve recorded by Hochster and D'Cruz.
\end{abstract}

\maketitle

\section{Introduction}

Moh constructed a family of height-two prime ideals in three-variable power-series rings whose minimal numbers of generators are unbounded in characteristic zero. He also asked whether these curves are set-theoretic complete intersections \cite[p.~733, Remark~1]{Moh1974}. The case $P_3$, with parametrization $(t^6+t^{31},t^8,t^{10})$, appears explicitly in Hochster's list \cite[Question~5]{Hochster2004}, Grifo's thesis \cite[Example~4.12]{Grifo2018}, and D'Cruz's survey \cite[Question~2.5]{DCruz2020}.

Moh subsequently studied the minimal generators and determinantal structure of his ideals \cite{Moh1979}. Mehta, Saha, and Sengupta gave a procedure for constructing explicit polynomial generators for the family in characteristic zero \cite{MehtaSahaSengupta2021}. For the parametrization above, Laura Gonz\'alez and Francesc Planas-Vilanova displayed four polynomial generators and proved that they minimally generate the formal prime in characteristic zero \cite[Discussion~4.6 and Theorem~4.7]{GonzalezPlanas2026}. The construction below starts from those generators and their relations.

The determinantal method is also classical. Valla showed that the maximal-minor ideal of a partly symmetric matrix is defined up to radical by two equations \cite[Theorem~2.1]{Valla1980}. Kleiman and Ulrich established the full ordinary-square containment in the proof of \cite[Lemma~2.5(3)]{KleimanUlrich1997}. The point here is an explicit invertible polynomial row operation that brings the matrix for this particular Moh curve into that setting. Its determinant is a nonzero constant, so it works over the whole polynomial ring.

Fix a field $k$ of characteristic zero, set $S=k[x,y,z]$, and write
\[
 \rho:S\longrightarrow k[t],\qquad
 (x,y,z)\longmapsto(t^6+t^{31},t^8,t^{10}),
 \qquad \mathfrak p=\ker\rho.
\]
The domain $S/\mathfrak p$ has transcendence degree one over $k$, so $\mathfrak p$ has height two. The arithmetic rank $\ara_S(\mathfrak p)$ is the least number of elements of $S$ generating an ideal with radical $\mathfrak p$.

\begin{theorem}\label{thm:main}
The polynomials $H_1,H_2$ defined in \eqref{eq:H} have integer coefficients and total degree $16$, and satisfy
\[
 \mathfrak p^2\subseteq(H_1,H_2)\subseteq\mathfrak p.
\]
Consequently,
\[
 \sqrt{(H_1,H_2)}=\mathfrak p,
 \qquad \ara_S(\mathfrak p)=2.
\]
\end{theorem}

Here $\mathfrak p^2$ denotes the ordinary ideal square. The same polynomials define the formal prime in $k[[x,y,z]]$ up to radical; see Corollary~\ref{cor:formal}.

The proof has three steps. First, the four known generators are shown to generate the polynomial kernel: an explicit inverse parameter identifies the localization away from $y=0$, and a nonzerodivisor argument recovers equality in $S$. Next, the row operation produces a symmetric maximal block without changing the ideal of maximal minors. Finally, a determinant and a bordered determinant give the displayed equations. Section~\ref{sec:symmetric} recalls the determinantal identity with a proof; Section~\ref{sec:moh} carries out the construction; and Section~\ref{sec:formal} passes to the completion. The result concerns the displayed parametrization; no assertion about the other members of Moh's family is needed.

\section{A symmetric maximal-minor lemma}\label{sec:symmetric}

Let $A$ be a commutative ring, let $M\in M_3(A)$ be symmetric, and let $r=(r_1,r_2,r_3)$. Set
\[
 B=\begin{pmatrix}M\\r\end{pmatrix},\qquad
 I=I_3(B),\qquad
 \Delta=\det M,\qquad C=\adj(M),
\]
and put
\[
 v=Cr^{\mathsf T}=(v_1,v_2,v_3)^{\mathsf T},
 \qquad \Gamma=rCr^{\mathsf T}=\sum_{a=1}^3 r_av_a.
\]
The matrix $C$ is symmetric. Expanding the maximal minors of $B$ along its last row gives
\[
 I=(\Delta,v_1,v_2,v_3).
\]

\begin{proposition}[Valla; Kleiman--Ulrich]\label{prop:symmetric}
With this notation,
\[
 I^2\subseteq(\Delta,\Gamma)\subseteq I.
\]
In particular, $\sqrt I=\sqrt{(\Delta,\Gamma)}$.
\end{proposition}

The two equations are those of Valla's construction \cite[Theorem~2.1]{Valla1980}; the full square containment is established in the proof of \cite[Lemma~2.5(3)]{KleimanUlrich1997}. More explicitly, take the symmetric matrix
\[
 \varphi=\begin{pmatrix}0&r\\r^{\mathsf T}&M\end{pmatrix}
\]
in the notation of that lemma. Its determinant is $-\Gamma$, its lower-right three-by-three minor is $\Delta$, and the maximal minors of its last three columns generate $I$. The containment follows from the identities alone, without the grade hypothesis used there to deduce self-linkage. The following proof recalls the identity needed here.

\begin{proof}
The expression $\Gamma=\sum_a r_av_a$ gives $(\Delta,\Gamma)\subseteq I$. For the other inclusion, expand
\begin{equation}\label{eq:minor-contraction}
 v_iv_j-\Gamma C_{ij}
 =\sum_{a,b=1}^3 r_ar_b
   \bigl(C_{ia}C_{jb}-C_{ij}C_{ab}\bigr).
\end{equation}
Since $C$ is symmetric, the parenthesis is
\[
 \det\begin{pmatrix}C_{ia}&C_{ij}\\C_{ba}&C_{bj}\end{pmatrix}.
\]
It is zero when the rows or columns repeat; otherwise it is, up to sign, a two-by-two minor of $C$. Jacobi's complementary-minor identity expresses every such minor as $\Delta$ times the complementary entry of $M$, up to sign \cite[p.~467]{Lawrence1964}. This is a polynomial identity over $\mathbb Z$: one may first prove it for the universal symmetric matrix after inverting its determinant, then use injectivity of that localization and specialize. It therefore holds over $A$, without a hypothesis on $\Delta$ or on the characteristic.

Equation~\eqref{eq:minor-contraction} now gives $v_iv_j\in(\Delta,\Gamma)$ for all $i,j$. Every generator of $I^2$ either contains a factor $\Delta$ or is one of these products, proving the assertion.
\end{proof}

After transposition, $B$ is partly symmetric in Valla's terminology. The second equation in his construction is the bordered determinant
\[
 \det\begin{pmatrix}M&r^{\mathsf T}\\r&0\end{pmatrix}=-\Gamma.
\]
Thus Proposition~\ref{prop:symmetric} uses the same two-generated ideal as \cite[Theorem~2.1]{Valla1980}. The argument above also works for a symmetric block of any size $n\geq2$: a two-by-two minor of its adjugate is its determinant times a complementary $(n-2)\times(n-2)$ minor. We use only the three-by-three case.

\section{Application to Moh's \texorpdfstring{$P_3$}{P3}}\label{sec:moh}

\subsection{The polynomial kernel}

Consider the polynomials of Gonz\'alez and Planas-Vilanova \cite[Theorem~4.7]{GonzalezPlanas2026}:
\begin{align*}
 f_1&=x^4-x^2y^4z^3-2xy^6z^2-4xyz-3y^3z^5+3y^3,\\
 f_2&=x^3y-xy^5z^3-3xz^2-2y^7z^2+2y^2z,\\
 f_3&=2x^3z-3x^2y^2-2xy^4z^4-y^6z^3+yz^2,\\
 f_4&=x^2yz-2xy^3-y^5z^4+z^3.
\end{align*}
Their cited result identifies the formal prime. For a global statement, one must also identify the kernel in the polynomial ring: a calculation in the completed local ring alone does not control the curve away from the origin. The next lemma establishes the required equality in $S$, using the fibre calculation and inverse parameter from Steps~3 and~6 of their proof, together with unmixedness and localization.

Set $J=(f_1,f_2,f_3,f_4)$ and define
\begin{equation}\label{eq:Phi}
 \Phi=\begin{pmatrix}
 0&2xy&-2z(2xy^3z^2+y^5z-1)\\
 z&-2x^2&y(4x^2yz^3+2xy^3z^2-3)\\
 -2y&-3z&-x+6y^2z^4\\
 3x&3y&0
 \end{pmatrix}.
\end{equation}
The first two columns are the coefficient vectors of the relations
\[
 zf_2-2yf_3+3xf_4=0,\qquad
 2xyf_1-2x^2f_2-3zf_3+3yf_4=0,
\]
which occur, up to nonzero scalar multiples, in the proof of \cite[Proposition~5.6]{GonzalezPlanas2026}. The third column completes this matrix to a presentation with the following maximal minors. Direct expansion gives
\begin{equation}\label{eq:Phi-minors}
 \det(\Phi_{\widehat i})=6(-1)^{i-1}f_i
 \qquad(1\leq i\leq4),
\end{equation}
where $\Phi_{\widehat i}$ is obtained by deleting row $i$. In particular, $I_3(\Phi)=J$.

\begin{lemma}\label{lem:global-kernel}
One has $J=\mathfrak p$. Hence $\Phi$ is a Hilbert--Burch matrix for $\mathfrak p$ over $S$.
\end{lemma}

\begin{proof}
\emph{Height and unmixedness.} Substitution into the parametrization gives $J\subseteq\mathfrak p$. Put $g=y^5-z^4$. The identity
\begin{equation}\label{eq:g-certificate}
 3g=y^2f_1-xyf_2-3zf_4
\end{equation}
shows that $g\in J$. The polynomials $g$ and $f_4$ have no common irreducible factor: such a factor would lie in $k[y,z]$ and divide the coefficient $yz$ of $x^2$ in $f_4$, whereas $g$ is relatively prime to $yz$. Thus $J$ has height at least two. Since $J\subseteq\mathfrak p$ and $\operatorname{ht}(\mathfrak p)=2$, we have $\operatorname{ht}(J)=2$.

Since $S$ is regular, $J$ has grade two. Equation~\eqref{eq:Phi-minors} and the Hilbert--Burch theorem give a free resolution of $S/J$ of length two \cite[Theorem~20.15, pp.~502--503]{Eisenbud1995}. If $\mathfrak q$ is an associated prime of $S/J$, then $(S/J)_{\mathfrak q}$ has depth zero. The Auslander--Buchsbaum formula over the regular local ring $S_{\mathfrak q}$ \cite[Tag~090U]{Stacks} therefore gives
\[
 \operatorname{ht}(\mathfrak q)
 =\operatorname{pd}_{S_{\mathfrak q}}(S/J)_{\mathfrak q}\leq2.
\]
Since $\mathfrak q$ contains $J$, its height is at least two. Thus all associated primes have height two. To pass from a localization back to $S$, it remains to show that $y$ is a nonzerodivisor.

\emph{The fibre at $y=0$.} Direct substitution gives
\[
 J+(y)=(y,x^4,xz^2,x^3z,z^3),
 \qquad \sqrt{J+(y)}=(x,y,z).
\]
No associated prime of $S/J$ can therefore contain $y$. It follows that $y$ is a nonzerodivisor on $S/J$, so
\begin{equation}\label{eq:y-injection}
 S/J\lhook\joinrel\longrightarrow(S/J)_y
\end{equation}
is injective.

\emph{The inverse parameter away from $y=0$.} We identify the localization explicitly. In
\[
 T=\bigl(S/(g,f_4)\bigr)_y,
\]
both $y$ and $z$ are units because $z^4=y^5$. Set $u=z/y$. Then $u^4=y$ and $u^5=z$, and $f_4=0$ becomes
\[
 0=u^9\bigl((x-u^3)^2-u^{31}\bigr).
\]
Consequently $w=(x-u^3)u^{-15}$ satisfies $w^2=u$ and is a unit. These identities give
\[
 x=w^6+w^{31},\qquad y=w^8,\qquad z=w^{10}.
\]
They exhibit mutually inverse homomorphisms
\[
 T\xrightarrow{\ \rho\ }k[t,t^{-1}],\qquad
 k[t,t^{-1}]\longrightarrow T,\quad t\longmapsto w.
\]
Indeed, $\rho(u)=t^2$ and $\rho(w)=t$. Hence $(g,f_4)S_y=\mathfrak pS_y$. Since $(g,f_4)\subseteq J\subseteq\mathfrak p$, we obtain $J S_y=\mathfrak pS_y$. The injection \eqref{eq:y-injection} then implies $J=\mathfrak p$.
\end{proof}

\subsection{The global symmetrization}

Proposition~\ref{prop:symmetric} applies once three rows of the presentation matrix form a symmetric block. An invertible row operation preserves the maximal-minor ideal, so the task is to find such an operation over $S$. The following matrix supplies it explicitly.

Set
\[
 \alpha=16x^3yz^2-10x^2y^3z+18y^2z^3,
 \qquad \beta=16x^2y^2z^2-10xy^4z,
\]
and define
\begin{equation}\label{eq:L}
 L=\begin{pmatrix}
 0&4&0&0\\
 -9&0&0&-\dfrac83x\\
 \alpha&\beta&6&6y^4z^2\\
 0&0&0&1
 \end{pmatrix}.
\end{equation}
Then
\[
 \det L=216,\qquad
 L\Phi=\begin{pmatrix}M\\r\end{pmatrix},\qquad
 r=(3x,3y,0),
\]
where
\begin{equation}\label{eq:M}
 M=\begin{pmatrix}
 4z&-8x^2&m_{13}\\
 -8x^2&-26xy&m_{23}\\
 m_{13}&m_{23}&m_{33}
 \end{pmatrix}
\end{equation}
and
\begin{align*}
 m_{13}&=4y(4x^2yz^3+2xy^3z^2-3),\\
 m_{23}&=18z(2xy^3z^2+y^5z-1),\\
 m_{33}&=32x^3yz^3-68x^2y^3z^2+30xy^5z-6x\\
       &\hspace{1em}-72xy^5z^6-36y^7z^5+72y^2z^4.
\end{align*}

The role of $\alpha$ and $\beta$ can be seen directly from the product $L\Phi$. The first two rows already have matching off-diagonal entries $-8x^2$. Matching the remaining entries in the third row requires
\[
 z\beta+18xy^4z^2-12y=m_{13},\qquad
 2xy\alpha-2x^2\beta+18y^5z^2-18z=m_{23},
\]
which the displayed choices satisfy. This proves the asserted symmetry without a localization.

Since $216$ is a unit in $k$, $L$ is invertible over $S$. The maximal minors of $L\Phi$ are $S$-linear combinations of those of $\Phi$; applying the same observation to $L^{-1}$ gives
\begin{equation}\label{eq:minor-ideal}
 I_3\!\begin{pmatrix}M\\r\end{pmatrix}=I_3(\Phi)=\mathfrak p.
\end{equation}

\subsection{The two equations}

The preceding construction gives the equations as $\det M$ and $r\adj(M)r^{\mathsf T}$. For ease of use, the next formulas express scalar multiples of these determinants as combinations of the four known generators.

Define
\begin{align*}
 \ell&=64x^4yz^2-40x^3y^3z+72xy^2z^3-81y^4z^2,\\
 \kappa&=-40x^3y^2z^2+25x^2y^4z+36y^3z^3,
\end{align*}
and set
\begin{equation}\label{eq:H}
 \begin{aligned}
 H_1&=-24xf_1+\ell f_3+81f_4,\\
 H_2&=-12yf_1+27xf_2+\kappa f_3.
 \end{aligned}
\end{equation}
Both polynomials belong to $\mathfrak p$ and have integer coefficients. Their total degrees are at most $16$. The monomial $x^5y^5z^6$ occurs in $H_1$ with coefficient $-128$, and $x^4y^6z^6$ occurs in $H_2$ with coefficient $80$, so both degrees equal $16$.

\begin{lemma}\label{lem:det}
For $\Delta=\det M$ and $\Gamma=r\adj(M)r^{\mathsf T}$, one has
\[
 \Delta=-16H_1,\qquad \Gamma=36H_2.
\]
\end{lemma}

\begin{proof}
Write $D_i=\det(\Phi_{\widehat i})=6(-1)^{i-1}f_i$. Expanding the determinant of the first three rows of $L\Phi$ gives
\begin{align*}
 \Delta&=64xD_1-\frac83\ell D_3+216D_4\\
       &=384xf_1-16\ell f_3-1296f_4=-16H_1.
\end{align*}
Let $D'_i$ denote the determinant obtained from $L\Phi$ by deleting row $i$. Expanding in the same way gives
\begin{align*}
 D'_1&=-54D_2-9\beta D_3=324f_2-54\beta f_3,\\
 D'_2&=24D_1-4\alpha D_3=144f_1-24\alpha f_3.
\end{align*}
If $v=\adj(M)r^{\mathsf T}$, the row-deletion signs give $D'_1=v_1$ and $D'_2=-v_2$. Therefore
\begin{align*}
 \Gamma&=3xD'_1-3yD'_2\\
       &=-432yf_1+972xf_2+(72y\alpha-162x\beta)f_3\\
       &=36H_2,
\end{align*}
where the last equality uses $72y\alpha-162x\beta=36\kappa$.
\end{proof}

\begin{proof}[Proof of Theorem~\ref{thm:main}]
Proposition~\ref{prop:symmetric} and \eqref{eq:minor-ideal} give
\[
 \mathfrak p^2\subseteq(\Delta,\Gamma)\subseteq\mathfrak p.
\]
Since $16$ and $36$ are units in $k$, Lemma~\ref{lem:det} identifies $(\Delta,\Gamma)$ with $(H_1,H_2)$. Taking radicals proves $\sqrt{(H_1,H_2)}=\mathfrak p$. Finally, Krull's height theorem gives $\ara_S(\mathfrak p)\geq\operatorname{ht}(\mathfrak p)=2$ \cite[Tag~0BBZ]{Stacks}, and the two equations give equality.
\end{proof}

For example, the square containment includes the identity
\[
 162f_3^2=-y(5x^3-2yz)H_1-x(8x^3+13yz)H_2.
\]

\section{The formal prime}\label{sec:formal}

The polynomial result yields the formal statement once the extended ideal is identified with the kernel of the power-series parametrization. The argument uses a finite extension followed by completion. Compare the normalization calculations of Gonz\'alez and Planas-Vilanova \cite[proof of Theorem~4.7, Steps~5--7]{GonzalezPlanas2026}.

Let $R=k[[x,y,z]]$ and let $P$ be the kernel of the continuous homomorphism
\[
 R\longrightarrow k[[t]],\qquad
 (x,y,z)\longmapsto(t^6+t^{31},t^8,t^{10}).
\]

\begin{corollary}\label{cor:formal}
One has $P=\mathfrak pR$ and
\[
 P^2\subseteq(H_1,H_2)R\subseteq P.
\]
In particular, $\sqrt{(H_1,H_2)R}=P$.
\end{corollary}

\begin{proof}
Put $B=S/\mathfrak p=k[t^6+t^{31},t^8,t^{10}]$ and $\mathfrak n=(x,y,z)B$. Since $t^8=y$, the ring $k[t]=B[t]$ is finite over $B$. Every prime of $k[t]$ lying above $\mathfrak n$ contains $t^8$, hence contains $t$ and equals $(t)$. Moreover $(t)\cap B=\mathfrak n$. Thus $(B\setminus\mathfrak n)^{-1}k[t]$ is local and equals $k[t]_{(t)}$, giving a finite injection
\[
 B_{\mathfrak n}\lhook\joinrel\longrightarrow k[t]_{(t)}.
\]
Completion with respect to $\mathfrak n$ preserves injections of finite modules over the Noetherian ring $B_{\mathfrak n}$ \cite[Tag~00MA]{Stacks}. In $k[t]_{(t)}$, the factor $1+t^{25}$ is a unit, so
\[
 \mathfrak n k[t]_{(t)}
 =(t^6(1+t^{25}),t^8,t^{10})=(t^6).
\]
The $\mathfrak n$-adic completion of this ring is therefore $k[[t]]$. Since the completion of $B_{\mathfrak n}$ is $R/\mathfrak pR$, we obtain an injection
\[
 R/\mathfrak pR\lhook\joinrel\longrightarrow k[[t]]
\]
induced by the parametrization. This proves $P=\mathfrak pR$. Extending the containments of Theorem~\ref{thm:main} to $R$ proves the assertions.
\end{proof}

\section*{Formal verification}

The kernel identifications, radical equalities, and arithmetic-rank conclusions for the polynomial and formal-power-series ideals are formalized in Lean~4 and registered in Palomar \cite{Gurung2026Formalization}. The registered statements use the polynomials $H_1,H_2$ of \eqref{eq:H} and hold over every field of characteristic zero. The citation concerns those statements; the proof of the ordinary-square containment is given in this note. The registered formalization is separate from the source attribution and historical discussion.

\section*{Use of AI tools}

Codex, ChatGPT, and Aristotle were used in developing the arguments, checking computations, and preparing the manuscript.

\begin{thebibliography}{99}

\bibitem{DCruz2020}
C.~D'Cruz,
\emph{Symbolic powers, set-theoretic complete intersection and certain invariants},
arXiv:2003.09101 (2020),
\url{https://arxiv.org/abs/2003.09101}.

\bibitem{Eisenbud1995}
D. Eisenbud, {\it Commutative algebra}, Graduate Texts in Mathematics, 150, Springer, New York, 1995; \href{https://mathscinet.ams.org/mathscinet/article?mr=1322960}{MR1322960}.

\bibitem{GonzalezPlanas2026}
Laura Gonz\'alez and Francesc Planas-Vilanova,
\emph{Prime ideals of Moh and the characteristic of the field},
J. Algebra Appl., published online 17 March 2026, article 2750170,
\url{https://doi.org/10.1142/S0219498827501702};
arXiv:2407.21692v2 (23 February 2026).

\bibitem{Grifo2018}
E.~Grifo,
\emph{Symbolic Powers and the Containment Problem},
Ph.D. thesis, University of Virginia, 2018,
\url{https://doi.org/10.18130/V3707WN5T}.

\bibitem{Gurung2026Formalization}
A.~Gurung,
\emph{Moh's $P_3$ curve is a set-theoretic complete intersection},
Lean~4 formalization, Palomar registry,
\href{https://palomar-registry.org/entry.html?id=PALOMAR-2026-09-07-000012&version=1}{PALOMAR-2026-09-07-000012},
version~1, 7 September 2026.

\bibitem{Hochster2004}
M.~Hochster,
\emph{Thirteen open questions in commutative algebra},
edited slides from a talk at Purdue University, July 2004,
\href{https://sites.lsa.umich.edu/hochster/wp-content/uploads/sites/1337/2024/10/Lip.text_.pdf}{author's online notes}.

\bibitem{KleimanUlrich1997}
S.~L. Kleiman and B. Ulrich, Gorenstein algebras, symmetric matrices, self-linked ideals, and symbolic powers, Trans. Amer. Math. Soc. {\bf 349} (1997), no.~12, 4973--5000; \href{https://mathscinet.ams.org/mathscinet/article?mr=1422609}{MR1422609}.
The cited numbering follows the preprint,
\href{https://arxiv.org/abs/alg-geom/9509005}{arXiv:alg-geom/9509005}.

\bibitem{Lawrence1964}
P.~A. Lawrence,
\emph{A generalization of Jacobi's theorem to hyperadjugates},
Canad. Math. Bull. \textbf{7} (1964), 467--469,
\url{https://doi.org/10.4153/CMB-1964-045-0}.

\bibitem{MehtaSahaSengupta2021}
R. Mehta, J. Saha and I. Sengupta, Moh's example of algebroid space curves, J. Symbolic Comput. {\bf 104} (2021), 168--182; \href{https://mathscinet.ams.org/mathscinet/article?mr=4180123}{MR4180123}.
Preprint: \href{https://arxiv.org/abs/1807.04909v4}{arXiv:1807.04909v4}.

\bibitem{Moh1974}
T.~T. Moh, On the unboundedness of generators of prime ideals in power series rings of three variables, J. Math. Soc. Japan {\bf 26} (1974), 722--734; \href{https://mathscinet.ams.org/mathscinet/article?mr=354680}{MR0354680}.

\bibitem{Moh1979}
T.~T. Moh, On generators of ideals, Proc. Amer. Math. Soc. {\bf 77} (1979), no.~3, 309--312; \href{https://mathscinet.ams.org/mathscinet/article?mr=545586}{MR0545586}.

\bibitem{Stacks}
The Stacks Project Authors,
\emph{The Stacks Project},
Tags \href{https://stacks.math.columbia.edu/tag/0BBZ}{0BBZ},
\href{https://stacks.math.columbia.edu/tag/00MA}{00MA},
and \href{https://stacks.math.columbia.edu/tag/090U}{090U},
\url{https://stacks.math.columbia.edu/}.

\bibitem{Valla1980}
G. Valla, On determinantal ideals which are set-theoretic complete intersections, Compositio Math. {\bf 42} (1980/81), no.~1, 3--11; \href{https://mathscinet.ams.org/mathscinet/article?mr=594479}{MR0594479}.

\end{thebibliography}


\end{document}
